\begin{table*}[!htbp]
\centering\small
\caption{Participant characteristics and recording devices.}
\label{tab:participants}
\setlength{\tabcolsep}{5pt}\renewcommand{\arraystretch}{1.14}
\begin{tabularx}{\textwidth}{@{}>{\raggedright\arraybackslash}Xrr@{}}
\toprule
Characteristic & ADHD study group & Control study group\\
\midrule
Participants & 41 & 55\\
Age, mean (SD), years & 11.51 (3.03) & 12.71 (3.88)\\
Age range, years & 7--18 & 6--18\\
Male & 33 & 27\\
Female & 8 & 28\\
\texttt{diagnosed=yes} & 41 & 13\\
\texttt{diagnosed=no} & 0 & 30\\
\texttt{diagnosed=undetermined} & 0 & 12\\
\midrule
Epoc X & 16 & 38\\
Epoc+ & 25 & 17\\
\bottomrule
\end{tabularx}
\par\smallskip\begin{minipage}{\textwidth}\footnotesize SD describes variation within group; the dataset README defines its gender codes as 1 = male and 2 = female. The available records do not explain the control-group entries with \texttt{diagnosed=yes}; neither field was recoded or treated as an independently re-adjudicated diagnosis.\end{minipage}
\end{table*}
